\documentclass[10pt]{article}
\usepackage[T1]{fontenc}
\usepackage[utf8]{inputenc}
\usepackage[french]{babel}
\usepackage{amsmath,amssymb,mathrsfs,mathpazo}
\usepackage{geometry}
\geometry{a4paper,margin=2cm}
\usepackage{graphicx,array,tabularx,tabularray,colortbl,multirow,multicol,booktabs,enumitem,pifont,ulem,fancybox,diagbox,cancel}
\usepackage{pstricks,pstricks-add,pst-plot,pst-tree,pst-node,pst-text,pst-func,pst-math,pst-bspline,pst-3dplot,pst-eucl,pst-solides3d}
\let\pscalculate\undefined
\usepackage{pst-bezier}
\usepackage{tikz,pgfplots,tkz-tab}
\usetikzlibrary{arrows,arrows.meta,patterns,calc,positioning,trees,decorations.markings,shapes.geometric,shapes.misc,shapes.arrows,intersections,decorations.pathmorphing,decorations.pathreplacing,backgrounds,shadows,babel}
\pgfplotsset{compat=1.18}
\usepgfplotslibrary{fillbetween}
\usepackage[np]{numprint}
\usepackage{listings,makecell,pas-tableur,twoopt}
\let\barre\undefined
\usepackage[pstricks]{tablvar}
\definecolor{darkgreen}{RGB}{0,100,0}
\providecommand{\red}{\color{red}}
\providecommand{\blue}{\color{blue}}
\providecommand{\green}{\color{green}}
\providecommand{\black}{\color{black}}
\providecommand{\white}{\color{white}}
\providecommand{\gray}{\color{gray}}
\providecommand{\cyan}{\color{cyan}}
\providecommand{\magenta}{\color{magenta}}
\providecommand{\yellow}{\color{yellow}}
\providecommand{\violet}{\color{violet}}
\providecommand{\orange}{\color{orange}}
\providecommand{\euro}{\texteuro}
\providecommand{\EUR}{\texteuro}
\providecommand{\ieme}{\textsuperscript{e}}
\providecommand{\degres}{\ensuremath{^\circ}}
\providecommand{\degree}{\ensuremath{^\circ}}
\providecommand{\vv}[1]{\overrightarrow{#1}}
\providecommand{\Coord}[1]{#1}
\providecommand{\decosix}{}
\providecommand{\starredbullet}{\ensuremath{\bullet}}
\pagestyle{empty}
\setlength{\parindent}{0pt}
\usepackage{adjustbox}\DeclareUnicodeCharacter{2236}{:}\DeclareUnicodeCharacter{2190}{\ensuremath{\leftarrow}}
\providecommand{\R}{\mathbb{R}}
\providecommand{\N}{\mathbb{N}}
\providecommand{\Z}{\mathbb{Z}}
\providecommand{\Q}{\mathbb{Q}}
\providecommand{\D}{\mathbb{D}}
\providecommand{\C}{\mathbb{C}}
\providecommand{\e}{{\mathrm{e}}}
\providecommand{\vect}[1]{\overrightarrow{#1}}
\providecommand{\vectt}[1]{\overrightarrow{\mathrm{#1}}}
\providecommand{\vvt}[1]{\overrightarrow{\mathrm{#1}}}
\providecommand{\np}[2][]{{#2~#1}}
\providecommand{\vv}[1]{{\overrightarrow{#1}}}
\renewcommand{\d}{{\mathrm{d}}}
\providecommand{\degres}{{^\circ}}
\providecommand{\degre}{{^\circ}}
\providecommand{\blue}{\color{blue}}
\providecommand{\red}{\color{red}}
\providecommand{\enskip}{\quad}
\providecommand{\llbracket}{\lbrack\!\lbrack}
\providecommand{\rrbracket}{\rbrack\!\rbrack}
\providecommand{\euro}{{\text{€}}}
\providecommand{\no}{{\text{n°}}}
\providecommand{\interval}[2]{{[#1;#2]}}
\providecommand{\intervFO}[2]{{[#1;#2[}}
\providecommand{\dotfill}{\ldots}

\providecommand{\ds}{\displaystyle}
\providecommand{\textding}[1]{#1}
\providecommand{\Oij}{\(\mathrm{O};\vec{\imath},\vec{\jmath}\)}
\providecommand{\Oijk}{\(\mathrm{O};\vec{\imath},\vec{\jmath},\vec{k}\)}
\providecommand{\Ouv}{\(\mathrm{O};\vec{u},\vec{v}\)}
\providecommand{\intervFF}[2]{\ensuremath{[#1;#2]}}
\renewcommand{\euro}{\ifmmode\text{\texteuro}\else\texteuro\fi}
\DeclareUnicodeCharacter{00B0}{\ifmmode^{\circ}\else\textdegree\fi}
\DeclareUnicodeCharacter{20AC}{\ifmmode\text{\texteuro}\else\texteuro\fi}
% Énoncés : Éducation nationale. Transcription et corrigés : APMEP, Denis Vergès et contributeurs. Découpage et classement : CoopMaths, Sébastien Lozano, Éric Elter et contributeurs. Adaptation de la présentation : Maths Explorer.
% CC BY-SA : https://coopmaths.fr/www/about/
% https://coopmaths.fr/alea/static/bac/2024/tex/bac_2024_05_sujet1_ameriquenord_1_cor.tex
\begin{document}
\noindent\textit{APMEP / CoopMaths. CC BY-SA. Présentation adaptée par Maths Explorer.}\par\medskip
\medskip















\textbf{Partie A}

\medskip











\begin{enumerate}
\item 
On dresse l'arbre pondéré de probabilités en utilisant les données de l'énoncé :

$p(R) = 0,07 \:; p_R(E) = 0,8$ et $p_{\overline{R}}(E) = 0,4$:

\begin{center}
\pstree[treemode=R,levelsep=3cm]{\TR{}}
{\pstree{\TR{$R~~$}\taput{0,07}}
	{\TR{$E$}\taput{0,8}
	\TR{$\overline{E}$}\tbput{\red 0,2}
	}
\pstree{\TR{$\overline{R}~~$}\tbput{\red 0,93}}
	{\TR{$E$}\taput{0,4}
	\TR{$\overline{E}$}\tbput{\red 0,6}
	}
}
\end{center}

On a $p(R \cap E) = p(R) \times p_{R}(E) = 0,07 \times 0,8 = 0,056$.
\item 
D'après la loi des probabilités totales :

$p(E) = p(R \cap E) + p\left(\overline{R} \cap E\right)$.

Or $p\left(\overline{R} \cap E\right) = p\left(\overline{R}\right) \times p_{\overline{R}}(E) = 0,93 \times 0,4 = 0,372$.

Donc $p(E) =0,056 + 0,372 = 0,428$.
\item 
On a $p_{E}(R) = \dfrac{p(E \cap R)}{p(E)} = \dfrac{p(R \cap E)}{p(E)} = \dfrac{0,056}{0,428} = \dfrac{56}{428} = \dfrac{14}{107} \approx \np{0,1308}$ soit 0,131 au millième près.
\end{enumerate}

\bigskip

\textbf{Partie B}

\medskip





\medskip

\begin{enumerate}
\item 
Les évènements étant indépendants et la probabilité d'obtenir un objet rare étant toujours égale à 0,07, la variable aléatoire $X$ suit une loi binomiale $\mathcal{B}(n,\:p)$ avec $n = 30$ et $p = 0,07$.

L'espérance mathématique est $E = n \times p = 30 \times 0,07 = 2,1$.
\item 

La calculatrice donne $P(X < 6) \approx \np{0,9838}$ soit 0,984 au millième près.
\item 
Quel que soit $k \in \N$, \: $p(X \geqslant k + 1) = 1 - p(X \leqslant k)$.
D'après la calculatrice :

$p(X \geqslant 1 + 1) = 1 - p(X \leqslant 1) \approx 0,631$, et $p(X \geqslant 2 + 1) = 1 - p(X \leqslant 2) \approx 0,351$, on a donc $k = 2$.

Dans le cadre du jeu ceci signifie que la probabilité d'obtenir au moins 2 objets rares est supérieure à $\dfrac12$.
\item 





Il faut donc trouver $X$ tel que $p(X \geqslant 1) \geqslant 0,95$ ou encore $1 - p(X = 0) \geqslant 0,95 \iff p(X = 0) \leqslant 1 - 0,95 \iff p(X = 0) \leqslant 0,05$.

Or $p(X = 0) = \binom{N}{0} \times 0,07^0 \times (1 - 0,07)^N = 0,93^N$.

Il faut donc résoudre l'inéquation :

$0,93^N \leqslant 0,05 \Rightarrow N \ln 0,93 \leqslant \ln 0,05$ par croissance de la fonction logarithme népérien  et enfin $N \geqslant \dfrac{\ln 0,05}{\ln 0,93}$, car $\ln 0,93 < 0$ et son inverse $\dfrac{1}{\ln 0,93}$ aussi.

D'après la calculatrice $\dfrac{\ln 0,05}{\ln 0,93} \approx 41,3$.

Conclusion Il faut que $N \geqslant 42$.
\end{enumerate}

\end{document}
